%% kapa: exam=2 round=30 year=2019 period=1 subject=appraisal_practice subject_ko=감정평가실무 seq=3 number=3 points=20
%% source: Q-Net 감정평가사 기출문제 (cst00309) artlSeq=5205595 file=raw/1008/5205595/감정평가실무.pdf
%% conversion: kordoc 4.24.51 -> tools/gen2.py (auto); review=proofread-sonnet; flags=none
\begin{longquestion}{【문제 3】}
감정평가사 甲은 부동산개발업자 乙로부터 개발계획의 타당성 검토를 의뢰받았다. 관련법규 및 이론을 참작하고 제시된 자료를 활용하여 개발계획의 타당성을 분석하시오. (20점) \par
< 자료 1 > 개발계획 \par
1. 부동산개발업자 乙은 K구 M동에 소재하는 노후된 상업용 부동산을 매수한 후 기존 건물을 철거하고 업무용건물을 신축하여 일정기간 임대한 후 처분할 계획임 \par
2. 매수 대상 부동산은 적정한 가격으로 매수 가능한 상태이며, 매수 대상 토지와 건물 중 기준시점 현재 건물의 잔존가치는 150,000,000원으로 판단됨 \par
3. 업무용건물의 신축공사 기간은 기준시점으로부터 1년이고, 1년 후 준공과 동시에 임대를 개시하며, 임대개시 5년 후 처분할 계획임 \par
4. 기준시점: 2019. 08. 01 \par
< 자료 2 > 매수 대상 부동산 \par
1. 소재지: K구 M동 300번지 \par
2. 토지: 대, 530 m\textsuperscript{2}, 상업용, 소로한면, 가장형, 평지 \par
3. 건물: 위 지상 벽돌조 슬라브지붕 지상 2층, 상업용, 건축연면적 630 m\textsuperscript{2} \par
4 용도지역: 준주거지역 \par % [원문 그대로] 번호 뒤 마침표 없음
< 자료 3 > 표준지 공시지가 (공시기준일: 2019. 01. 01.) \par
{\small\centering
\begin{tabular}{|c|c|c|c|c|c|c|c|c|}
\hline
\makecell{기호} & \makecell{소재지} & \makecell{면적(m\textsuperscript{2})} & \makecell{지목} & \makecell{이용\newline{}상황} & \makecell{용도\newline{}지역} & \makecell{도로\newline{}접면} & \makecell{형상\newline{}지세} & \makecell{공시지가\newline{}(원/m\textsuperscript{2})} \\
\hline
\makecell{1} & \makecell{K구 M동 293} & \makecell{400} & \makecell{대} & \makecell{상업용} & \makecell{준주거지역} & \makecell{세로가} & \makecell{가장형\newline{}평지} & \makecell{1,760,000} \\
\hline
\makecell{2} & \makecell{K구 M동 307} & \makecell{520} & \makecell{대} & \makecell{상업용} & \makecell{준주거지역} & \makecell{소로한면} & \makecell{가장형\newline{}평지} & \makecell{1,870,000} \\
\hline
\end{tabular}\par}

< 자료 4 > 거래사례 자료 \par
1. 소재지: K구 M동 315번지 \par
2. 토지: 대, 490m\textsuperscript{2}, 상업나지, 소로한면, 가장형, 평지 \par
3. 건물: 없음(토지만의 거래사례임) \par
4. 용도지역: 준주거지역 \par
5. 거래금액: 1,150,000,000원 \par
6. 거래시점: 2019. 05. 01 \par
7. 기타사항: 대상지역의 거래관행은 거래시점에 매매대금을 모두 일시에 지급하는 것이나, 본건 거래사례의 경우 매매대금의 70 \%를 거래시점에 지급하고 나머지 30 \%는 1년 후에 지급하는 조건인 것으로 확인되었으며, 그 외의 거래내역은 정상적임 \par
< 자료 5 > 시점수정, 지역요인, 개별요인, 그 밖의 요인 자료 \par
1. 시점수정치(지가변동률) \par
1) 2019. 01. 01∼2019. 08. 01: 1.01752 \par
2) 2019. 05. 01∼2019. 08. 01: 1.00697 \par
2. 지역요인 비교치: 대상 부동산, 표준지, 거래사례 모두 인근지역 내에 소재하여 지역요인은 동일함 \par
3. 개별요인 비교치: 개별요인 비교항목을 검토한 결과 대상 토지는 표준지 기호 1 대비 5 \% 우세하고, 표준지 기호 2 대비 1 \% 열세하며, 거래사례 대비 2 \% 우세한 것으로 격차율이 산정되었음 \par
4. 그 밖의 요인 보정치: 1.25를 적용함 \par
< 자료 6 > 건물 신축 관련 자료 \par
1. 구조 및 용도 등: 철근콘크리트조 슬라브지붕, 지하 2층ㆍ지상 7층, 업무용 \par
2. 면적: 건축면적 300 m\textsuperscript{2}, 건축연면적 2,700 m\textsuperscript{2} \par
3. 신축공사기간: 기준시점에 착공하여 1년 후 준공 \par
4. 건축공사비: 900,000 원/m\textsuperscript{2}(건축공사비에는 기존 건물의 철거비 및 기타 제반 부대비용 등이 포함되어 있으며, 착공시점에 30 \%, 준공시점에 70 \%를 지급함) \par
< 자료 7 > 임대수익 관련 자료 \par
1. 임대가능면적: 건축연면적의 70 \%임 \par
2. 순영업소득: 연간 순영업소득은 145,000 원/m\textsuperscript{2}이며, 보유기간 동안 변동 없이 유지될 것으로 판단됨 \par
< 자료 8 > 대출조건 및 기타 자료 \par
1. 대출조건: 대출비율 60 \%, 이자율 7 \%, 만기 20년임(매월 원리금 균등상환) \par
2. 자기자본수익률: 10 \% \par
3. 할인율: 8 \% \par
4. 임대개시 5년 후 처분할 계획이며, 부동산가치는 임대개시(준공) 이후 매년 2 \% 상승할 것으로 판단됨 \par
5. 수익환원은 직접환원법에 의함 \par
6. 환원율 계산시 소수점 이하 다섯째자리에서 반올림하여 소수점 이하 넷째자리까지 표시함 \par
7. 각 단계의 가격산정시 천원미만은 절사함 \par
\end{longquestion}
